\section{Einleitung}
\label{Hauptkapitel:Einleitung}
Die rasante Entwicklung von Large Language Models (LLMs) hat in den letzten Jahren beeindruckende Fortschritte in der automatisierten Verarbeitung natürlicher Sprache hervorgebracht. Weil sie in der Lage sind, komplexe Fragestellungen zu erfassen und sinnvolle Antworten darauf zu formulieren, finden fortschrittliche Modelle immer öfter Einsatz in wissensbasierten Anwendungen und Expertensystemen.  Trotz dieser Fortschritte gibt es ein erhebliches Problem:  Da die Inhalte von LLMs probabilistisch erstellt werden, können sie sogenannte Halluzinationen produzieren: Antworten, die zwar glaubwürdig wirken, aber nicht auf verifizierbaren Informationen beruhen.

In thematisch fokussierten Expertensystemen, in denen die Korrektheit und Nachvollziehbarkeit von Informationen entscheidend sind, stellt dieses Verhalten eine erhebliche Herausforderung dar. Klassische Expertensysteme basieren zwar auf strukturierten Wissensbasen und liefern nachvollziehbare Ergebnisse, sind jedoch häufig unflexibel und schwer skalierbar. Im Gegensatz dazu zeichnen sich LLMs durch eine hohe Flexibilität aus, allerdings auf Kosten der Verlässlichkeit.

Das Konzept der Retrieval-Augmented Generation (RAG) stellt einen vielversprechenden Ansatz dar, um die Stärken beider Ansätze zu kombinieren. Dabei werden externe Wissensquellen in den Generierungsprozess integriert, indem relevante Dokumente zunächst mithilfe von Information-Retrieval-Methoden identifiziert und anschließend als Kontext für die Antwortgenerierung verwendet werden. Moderne Suchsysteme wie Elasticsearch ermöglichen dabei eine sogenannte Hybrid Search, die sowohl lexikalische Verfahren (z.B. BM25) als auch semantische Vektorsuche kombiniert. Dadurch können sowohl exakte Wortübereinstimmungen als auch semantische Ähnlichkeiten berücksichtigt werden.

In diesem Zusammenhang besteht die zentrale Herausforderung darin, ein System zu entwickeln, das relevante Informationen zuverlässig identifiziert, die Antwortgenerierung strikt auf diese Informationen beschränkt und gleichzeitig transparente sowie nachvollziehbare Ergebnisse liefert. Darüber hinaus muss das System erkennen, wenn die vorhandenen Informationen nicht ausreichen, um eine Anfrage zu beantworten, und in solchen Fällen bewusst auf eine Antwort verzichten.

Vor diesem Hintergrund untersucht diese Arbeit die Nutzbarkeit von Large Language Models in einem thematisch fokussierten Expertensystem in Verbindung mit hybriden Retrieval-Methoden. Zu diesem Zweck wird ein System entwickelt, das Elasticsearch als Vektordatenbank nutzt, eine Kombination aus BM25-basierter Volltextsuche und semantischer Vektorsuche realisiert und die gewonnenen Ergebnisse gemeinsam mit einem LLM zur Beantwortung von Nutzeranfragen verwendet. Im Fokus stehen dabei die Qualität der generierten Antworten, die Minimierung von Halluzinationen sowie die Gewährleistung der Nachvollziehbarkeit durch die Angabe der zugrunde liegenden Textpassagen.

Die Struktur der Arbeit ist wie folgt: In Kapitel 2 werden die theoretischen Grundlagen zu Large Language Models, Information Retrieval und Hybrid Search erläutert. Kapitel 3 gibt einen Überblick über den aktuellen Stand der Forschung. In Kapitel 4 wird die Systemarchitektur vorgestellt, während Kapitel 5 die Implementierung beschreibt. Kapitel 6 widmet sich der Evaluation des Systems, gefolgt von einer kritischen Diskussion der Ergebnisse in Kapitel 7. Abschließend werden in Kapitel 8 die zentralen Erkenntnisse zusammengefasst und ein Ausblick auf zukünftige Entwicklungen gegeben.
\clearpage
